%#%%\FINALFIX{\newpage}%
\subsection{nml\_timestep}
\label{SEC:nml_timestep}
\inp*[nmlTimestep]{}[nml\_timestep][]

This \hyperref[SEC:User_Input]{namelist block}
is used to specify time stepping features.

\begin{nmlblocklst}[
,label={LST:nml_timestep}
,caption={Namelist nml\_timestep.}
]
&nml_timestep
  Runge_Kutta_Scheme        =    113
  Num_Timesteps_Min         =    100
  Num_Timesteps             =  25000
  Final_Time                =  -10.0  ! non-dim time, t*
  Timestep_Type             =      1
  Timestep                  = 1.0e-6  ! t[s] * ( a[m/s] / L=1.0[m])
  CFL                       =    1.0
  CFL_Ramping_Beg           =    0.1
  CFL_Ramping_End           =    1.0
  CFL_Ramping_Iterations    =      1
  CFL_Ramping_Start_Iter    =      0
  Convergence_Order_Abs_Res =  -16.0
  Convergence_Order_Max_Res =  -16.0
/
\end{nmlblocklst}

%==========================================================================

\begin{nmldescitem}
Runge_Kutta_Scheme = 113
\end{nmldescitem}
\inp*[nmlTimestep::RungeKuttaScheme]{}[Runge\_Kutta\_Scheme][]

\begin{nmldescbody}
Specifies which Runge-Kutta (RK) method to use.
%
See \sref{SEC:Time:Methods} for additional details.
%
Speedup is relative to the 3-stage/\nth{3}-order SSP ERK method.

\begin{longtable}{ r @{ = } r @{/} p{3.55in} c }
  \toprule
  \multicolumn{1}{l@{ \phantom{=} }}{Option} &
  \multicolumn{2}{@{}l}{Description} &
  Speedup \\
  \midrule%\\[-0.75em]
  \endhead
%\multicolumn{4}{l}{}\\[-0.75em]
         \multicolumn{1}{r@{}}{}
       & \multicolumn{2}{@{}l}{\textbf{Recommended options.}} & \\
   113 &  11-stage & \nth{3}-order VOERK\cite{AlJahdali_2022} & 1.33x \\
   154 &  15-stage & \nth{4}-order VOERK\cite{AlJahdali_2022} & 1.33x \\
     4 &   5-stage & \nth{4}-order SSP ERK\cite{Spiteri_2002} & 1.22x \\
   165 &  16-stage & \nth{5}-order VOERK\cite{AlJahdali_2022} & 1.18x \\
   943 &   9-stage & \nth{4}(\nth{3})-order
                     EPERK FSAL\cite{Ranocha_2022}            & 1.22x \\
   532 &   5-stage & \nth{3}(\nth{2})-order
                     EPERK FSAL\cite{Ranocha_2022}            & 1.09x
\\[0.5em]
\pagebreak[3]
%
         \multicolumn{1}{r@{}}{}
       & \multicolumn{2}{@{}l}{\textbf{Other options that work well.}}   & \\
    -2 &   2-stage & \nth{2}-order Classic RK                       & n/a   \\
    -3 &   3-stage & \nth{3}-order Classic RK                       & n/a   \\
    -4 &   4-stage & \nth{4}-order Classic RK                       & n/a   \\
     2 &   2-stage & \nth{2}-order SSP ERK\cite{Gottlieb_1998}      & n/a   \\
     3 &   3-stage & \nth{3}-order SSP ERK\cite{Gottlieb_1998}      & 1.00x \\
     5 &   5-stage & \nth{4}-order Carpenter-Kennedy low storage
                                   RK\cite{Carpenter_1994}          & 0.99x \\
    53 &   5-stage & \nth{3}-order VOERK\cite{AlJahdali_2022}       & 1.08x \\
    64 &   6-stage & \nth{4}-order VOERK\cite{AlJahdali_2022}       & 1.02x \\
    85 &   8-stage & \nth{5}-order VOERK\cite{AlJahdali_2022}       & 1.05x \\
  -943 &   9-stage & \nth{4}(\nth{3})-order
                     EPERK non-FSAL\cite{Ranocha_2022}              & n/a   \\
  -532 &   5-stage & \nth{3}(\nth{2})-order
                     EPERK non-FSAL\cite{Ranocha_2022}              & n/a
\\[0.5em]
\pagebreak[3]
%
         \multicolumn{1}{r@{}}{}
       & \multicolumn{2}{@{}l}{\textbf{Options that do not perform as well.}}&\\
  $\pm$1 &   1-stage & \nth{1}-order Forward Euler method       & n/a   \\
     -53 &   5-stage & \nth{3}-order AOERK\cite{AlJahdali_2022} & n/a \\
    -113 &  11-stage & \nth{3}-order AOERK\cite{AlJahdali_2022} & n/a \\
     -64 &   6-stage & \nth{4}-order AOERK\cite{AlJahdali_2022} & n/a \\
    -154 &  15-stage & \nth{4}-order AOERK\cite{AlJahdali_2022} & n/a \\
     -85 &   8-stage & \nth{5}-order AOERK\cite{AlJahdali_2022} & n/a \\
    -165 &  16-stage & \nth{5}-order AOERK\cite{AlJahdali_2022} & 0.98x \\
      54 &   6-stage & \nth{5}(\nth{4})-order
                       Dormand-Prince EPERK\cite{Dormand_1980}  & 0.64x \\
     853 &  12-stage & \nth{8}(\nth{5})-order
                       Dormand-Prince EPERK\cite{Hairer_1993}   & 0.61x \\
    1054 &  10-stage & \nth{5}(\nth{4})-order
                       EPERK FSAL\cite{Ranocha_2022}            & 0.85x \\
   -1054 &  10-stage & \nth{5}(\nth{4})-order
                     EPERK non-FSAL\cite{Ranocha_2022}        & n/a \\
  \bottomrule
\end{longtable}

\noindent
Acronyms used in the description above:
\\
\begin{tabular}{ @{\qquad} l @{ : } p{0.75\linewidth} @{} }
  SSP   & Strong stability preserving \\
  VOERK & Vortex-optimized explicit Runge-Kutta (ERK) \\
  AOERK & Advection-optimized ERK \\
  FSAL  & First-same-as-last \\
  EPERK & Embedded-pair ERK with adaptive time stepping
\end{tabular}
\\

\NOTE
EPERK methods do not work with limiting or filtering.

\NOTE
The value in parenthesis for EPERK methods is the order of the embedded method.
%
The difference between the actual method and embedded method represents the
time-error from the time step.
%
This is used by the adaptive time stepping proportional–integral–derivative
(PID) controller to compute the size of the next time step.
\end{nmldescbody}

%==========================================================================

\begin{nmldescitem}
Num_Timesteps_Min = 100
\end{nmldescitem}
\inp*[nmlTimestep::NumTimestepsMin]{}[Num\_Timesteps\_Min][]

\begin{nmldescbody}
The minimum number of time steps to complete before checking whether any of
the convergence criteria have been met.

\NOTE When initializing a simulation to a uniform flow, the very
small changes that occur during the first few time steps can trigger a false
positive convergence check if this is set to a small number,
\forexample $\leq5$.
\end{nmldescbody}

%==========================================================================

\begin{nmldescitem}
Num_Timesteps = 25000
\end{nmldescitem}
\inp*[nmlTimestep::NumTimesteps]{}[Num\_Timesteps][]

\begin{nmldescbody}
Number of time steps for the simulation to run, unless
\inp[nmlTimestep::FinalTime]{Final\_Time}
is reached before completing all time steps.
\end{nmldescbody}

%==========================================================================

\begin{nmldescitem}
Final_Time = -10.0  ! non-dim time, t*
\end{nmldescitem}
\inp*[nmlTimestep::FinalTime]{}[Final\_Time][]

\begin{nmldescbody}
Stops the simulation when the time within the simulation reaches this
final time value.
%
This value is specified in nondimensional units
(\key[key:tstar]{\tstar}[t *@\tstar]), \eref{EQN:tstar}, based on the input
provided for
\inp[nmlOutput::TimeScalingFactor]{Time\_Scaling\_Factor}.
%
For the isentropic vortex problem
(\inp[nmlInitialization::Initialization]{Initialization}=5 or 55), a
negative value specifies the total number of periods to run the simulation.
%
For all other cases, a negative value disables the use of
\inptext{Final\_Time} to stop a simulation.

\NOTE
The nondimensional simulation time
(\key[key:tstar]{\tstar}[t *@\tstar]), \eref{EQN:tstar},
is written to standard output with the convergence information
and to the residual file.
\end{nmldescbody}

%==========================================================================

\begin{nmldescitem}
Timestep_Type = 1
\end{nmldescitem}
\inp*[nmlTimestep::TimestepType]{}[Timestep\_Type][]

\nopagebreak
\begin{nmldescbody}
This determines the method used to compute the time step size and the type of
time stepping to use for the simulation (unless using one of the EPERK methods
for adaptive time stepping, see note below).
%
\\
\begin{tabular}{ @{\qquad} r @{ = } p{0.80\linewidth} @{} }
  -2 & Local time step from the inviscid eigenvalue. (Experimental) \\
  -1 & Local time step from the spectral radius. \\
   0 & Global time step from user specified time step. \\
   1 & Global time step from the spectral radius. \\
   2 & Global time step from the inviscid eigenvalue.  (Experimental)
\end{tabular}
 
\begin{InputList}[font=\mdseries]
  \item[Local time stepping:]
    At the beginning of each solver iteration, each grid cell computes its own
    unique time step size that is used to update the solution within the cell at
    the end of each Runge-Kutta stage.
    %
    Since all cells are updated with different time step sizes, this can only be
    used for steady-state simulations and cannot be used for time-accurate
    simulations (\forexample LES) which require all cells to be synchronized in
    time.
  \item[Global time stepping:]
    At the beginning of each solver iteration, the time step for each cell is
    computed in the same way as with local time stepping.
    %
    However, the minimum time step size across all grid cells is taken to be the
    global time step size that all cells use to update their solutions in time.
    %
    Global time stepping is required for time-accurate simulation like LES, but
    is generally inefficient for steady-state problems.
  \item[Inviscid eigenvalue:]
    The maximum inviscid eigenvalue is used to compute a time step size at
    each solution point within a cell, and the minimum of these values is
    taken to be the time step for the cell.
  \item[Spectral radius:]
    The time step within each cell is taken as the minimum value of the
    spectral radius evaluated at every face\slash flux point of the cell.
  \item[User specified time step:]
    A constant global time step is obtained from the user input for
    \inp[nmlTimestep::Timestep]{Timestep}.
    %
    \NOTE The input \inptext{Timestep} is specified as a nondimensional value.
    It is the physical time step nondimensionalized by the reference speed of
    sound and a length of 1 meter.
%
%   A constant global time step is computed from the inputs
%   $\inp[nmlTimestep::CFL]{CFL} \times
%    \inp[nmlTimestep::Timestep]{Timestep}$.
%   \newline
%   %
%   For this option, \inptext{CFL}=1.0 unless CFL ramping is being used.
%   See \inp[nmlTimestep::Timestep]{Timestep}.
\end{InputList}

\vspace{1.0ex}\noindent\NOTE
For each of the methods above,
once the time step is determined, the effective $\Delta t$ will be computed as
\begin{equation}
  \Delta t = CFL \times \text{time step}
\end{equation}
where $CFL$ is obtained from
  \inp[nmlTimestep::CFL]{CFL}
or from
  \inp[nmlTimestep::CFLRampingBeg]{CFL\_Ramping\_Beg}/%
  \inp[nmlTimestep::CFLRampingEnd]{End}[CFL\_Ramping\_End]
if CFL ramping is being used.
%
If \inptext{Timestep\_Type}=0 (user specified time step), then
\inp[nmlTimestep::CFL]{CFL} will be set to 1.0,
however ramping may still be used.
See \inp[nmlTimestep::Timestep]{Timestep}.
\\ % Adds an extra blank line.

\noindent\NOTE
If using one of the EPERK methods for adaptive time stepping, the CFL column
written to \file{stdout} for each iteration will contain an approximate CFL
value, $\text{CFL}_\textsc{eperk}$, for the adaptive time step size relative
to the time step size that would have been computed for one of the CFL-based
ERK methods, \thatis
%
\[
  \text{CFL}_\textsc{eperk} \approx
                   \frac{\Delta t_\textsc{eperk}}{\Delta t_\textsc{erk}}
\]
%
where $\Delta t_\textsc{eperk}$ is the time step size computed by the PID
controller for the EPERK method, and $\Delta t_\textsc{erk}$ is the time step
size that would have been computed using the method specified by
\inptext{Timestep\_Type} for a CFL-based ERK method with $\text{CFL}=1.0$.

\devNotes{
For \inptext{Timestep\_Type=0} (user specified time step), have the
code force CFL=1.0 rather than just ignore it.
}
\end{nmldescbody}

%==========================================================================

\begin{nmldescitem}
Timestep = 1.0e-6  ! t[s] * (a[m/s] / L=1.0[m])
\end{nmldescitem}
\inp*[nmlTimestep::Timestep]{}[Timestep][]

\begin{nmldescbody}
Specifies the size of the time step that is used when
\inp[nmlTimestep::TimestepType]{Timestep\_Type}=0
(constant global time step).
%\thatis $\inptext{Timestep\_Type}=0$.
%
The value specified by this input is nondimensionalized by
$L/a_{\infty}$, the reference length (1~m) and 
speed of sound.
Thus, the dimensional time step size will be
\inp[nmlTimestep::Timestep]{Timestep}$ * L / a_\infty$ seconds.
 
When a simulation is running, the nondimensional time step size is
the second column written to stdout, so the value of this input is
identical to that output under the conditions
\inp[nmlTimestep::TimestepType]{Timestep\_Type}=0 and
\inp[nmlTimestep::CFL]{CFL}=1.0.
%
%\CHANGECODE
When using \inp[nmlTimestep::TimestepType]{Timestep\_Type}=0
(constant global time step), the value used for
\inptext{CFL} is $1.0$ unless CFL ramping is being used.
%
If CFL ramping is being used, the ramping parameters will be used
to evaluate the value of \inptext{CFL} to be used for each time step,
but \inptext{CFL} will revert to a value of $1.0$ once the
ramping completes.

\NOTE
This nondimensional time step is written to standard output
with the convergence information.
%
The \inptext{Timestep} is nondimensionalized by
the reference length (1~m) and speed of sound, whereas
the solution time that is written to standard output and
\inp[nmlTimestep::FinalTime]{Final\_Time} are both nondimensionalized by
\inp[nmlOutput::TimeScalingFactor]{Time\_Scaling\_Factor}
via \eref{EQN:tstar}.

\end{nmldescbody}

%==========================================================================

\begin{nmldescitem}
CFL = 1.0
\end{nmldescitem}
\inp*[nmlTimestep::CFL]{}[CFL][]

\begin{nmldescbody}
Analogous to the CFL condition, this is used to scale the time step size to
get a maximum stable time step. The maximum value this can be while also
keeping a simulation stable is very dependent on the type of simulation,
flow conditions, and input settings.  Experimentation is generally
required to find a good value.  A lower value is more likely to keep a
simulation stable, but will require more time steps to reach either
convergence or a given simulation time. Typical values are
$0.1\leq\inptext{CFL}\leq2$, but past experience has found simulations that
were stable with values up to $10.0$ for certain Runge-Kutta schemes.

%\NOTE
If CFL ramping is actively being used, that computed value
will be applied.
Once CFL ramping has completed, the value of this input will be used.

\devNotes{
CHANGECODE:
This value will automatically be set to $1.0$ if using
\inp[nmlTimestep::TimestepType]{Timestep\_Type}=0
(constant global time step), unless CFL ramping is active.
}

\end{nmldescbody}

%==========================================================================

\begin{nmldescitem}
CFL_Ramping_Beg = 0.1
CFL_Ramping_End = 1.0
\end{nmldescitem}
\inp*[nmlTimestep::CFLRampingBeg]{}[CFL\_Ramping\_Beg][]
\inp*[nmlTimestep::CFLRampingEnd]{}[CFL\_Ramping\_End][]

\begin{nmldescbody}
Beginning and ending CFL values to use with the linear CFL ramping that is
active when
$\inp[nmlTimestep::CFLRampingIterations]{CFL\_Ramping\_Iterations}\geq2$.

See also \sref{SEC:Files:Interrupt:Adjust_CFL} for adjusting the CFL
while the solver is running.
\end{nmldescbody}

%==========================================================================

\begin{nmldescitem}
CFL_Ramping_Iterations = 1
\end{nmldescitem}
\inp*[nmlTimestep::CFLRampingIterations]{}[CFL\_Ramping\_Iterations][]

\begin{nmldescbody}
The number of time steps used to linearly ramp the CFL value,
starting at \inptext{CFL\_Ramping\_Beg} and
ending   at \inptext{CFL\_Ramping\_End}.
%
$\inptext{CFL\_Ramping\_Iterations}\geq2$ activates CFL ramping.

\NOTE
A new restart file will be written just before the
ramping schedule begins in case it is too aggressive and causes the
simulation to fail.

See also \sref{SEC:Files:Interrupt:Adjust_CFL} for adjusting the CFL
while the solver is running.
\end{nmldescbody}

%==========================================================================

\begin{nmldescitem}
CFL_Ramping_Start_Iter = 0
\end{nmldescitem}
\inp*[nmlTimestep::CFLRampingStartIter]{}[CFL\_Ramping\_Start\_Iter][]

\begin{nmldescbody}
Specifies the time step iteration at which to begin the CFL ramping schedule
specified by
\inp[nmlTimestep::CFLRampingBeg]{CFL\_Ramping\_Beg},
\inp[nmlTimestep::CFLRampingEnd]{CFL\_Ramping\_End}, and
\inp[nmlTimestep::CFLRampingIterations]{CFL\_Ramping\_Iterations}.
%
This input is primarily used with the \inptext{ADJUST\_CFL} control file
(\sref{SEC:Files:Interrupt:Adjust_CFL})
that can be used to modify the CFL ramping schedule of a running simulation.

\NOTE
A new restart file will be written just before the
ramping schedule begins in case it is too aggressive and causes the
simulation to fail.


See also \sref{SEC:Files:Interrupt:Adjust_CFL} for adjusting the CFL
while the solver is running.
\end{nmldescbody}

%==========================================================================

\begin{nmldescitem}
Convergence_Order_Abs_Res = -16.0
\end{nmldescitem}
\inp*[nmlTimestep::ConvergenceOrderAbsRes]{}[Convergence\_Order\_Abs\_Res][]

\begin{nmldescbody}
This is the convergence goal for the orders of magnitude of the absolute
residual.
\end{nmldescbody}

%==========================================================================

\begin{nmldescitem}
Convergence_Order_Max_Res = -16.0
\end{nmldescitem}
\inp*[nmlTimestep::ConvergenceOrderMaxRes]{}[Convergence\_Order\_Max\_Res][]

\begin{nmldescbody}
This is the convergence goal for the orders of magnitude reduction of the
residual relative to the maximum value of the residual during the
simulation.
\end{nmldescbody}

%==========================================================================

